\section{Introduction, claims, and the status of each claim}\label{sec:intro}

SpherePOP is a notation in which a computation is a finite tree of nested circles. A circle carries an operator and contains operands, which are numbers or further circles. The single evaluation act is a pop: a circle all of whose operands are values is replaced by the value its operator computes from them. The introductory text on which this essay builds develops the pure fragment in detail, and an executable companion verifies its basic facts. The question here is deliberately narrower and harder than the question of whether the notation is pleasant to use. It is the question of what the notation can compute, and what has to be added to it before it can compute everything a Turing machine can.

The pure fragment cannot. The reason is not subtle, and it is worth stating at the outset because the whole essay is organized around it. In the pure fragment every pop removes one circle and never creates one. A tree with $n$ circles therefore reaches a value in exactly $n$ pops, whatever order the pops are performed in, and a calculus in which every computation has a length fixed in advance by the size of its input cannot express an unbounded search. Universality requires that some rule be able to make a term larger, and in particular that some rule be able to duplicate a piece of the term, so that the same sub-computation can be demanded many times. The extension studied below is the smallest one that does this: a single binary \emph{application circle} $\App(f,x)$ and two contraction rules. Reading $\App$ as juxtaposition and associating to the left,
\[
\Kc\,x\,y \;\red\; x,
\qquad
\Sc\,x\,y\,z \;\red\; x\,z\,(y\,z).
\]
The first rule deletes the circles of $y$. The second contains $z$ twice on its right side and once on its left, which is the duplication. These are the combinators $\Sc$ and $\Kc$ of Sch\"onfinkel and of Curry \cite{schonfinkel,curry-feys}, and the calculus they generate is combinatory logic. The choice of $\Sc$ and $\Kc$ as the foundation, rather than lambda abstraction, is a choice about where to place the burden of proof. Combinatory logic has no bound variables, so there is no substitution to define, no capture to avoid, and no question about reduction under a binder. Every circle in a term is either an atom or an application, and every contraction is a local rewrite of a pattern whose shape is fixed. Lambda binding is then treated afterwards, as a source language that is translated into circles, so that its relationship to the circle calculus becomes a theorem about a translation and not an assumption about a primitive.

\subsection{What this changes in the book}

One departure from the introductory text must be flagged, because it is the point at which a careful reader would otherwise object. In the pure fragment an operator is applied only to values: a circle is popped when its operands are numbers, and evaluation proceeds from the leaves inward. The application circle as used here accepts arbitrary terms as operands. In $\Kc\,x\,y$ the operand $y$ may be a circle that has not been popped and, if it were popped first, might never finish. The rule $\Kc\,x\,y\red x$ discards $y$ unexamined. This is the difference between a call-by-value discipline, in which operands are reduced before the operator is applied, and a call-by-name discipline, in which they are not. It is not an incidental detail. It is precisely what makes the calculus sensitive to the order in which circles are popped, and Section~\ref{sec:nontermination} shows that the same term can have a value under one order and none under another. The pure fragment is the sub-calculus in which this sensitivity is absent, and Section~\ref{sec:pure} records exactly why.

\subsection{Claims}

The essay makes the following claims, and it keeps them apart because they are logically independent and are established by different means.

\begin{enumerate}[label=(\Alph*)]
\item \emph{Failure of the size measure under duplication} (Section~\ref{sec:size}). Contracting a $\Kc$ redex changes the number of circles by exactly $-(2+\size{y})$ and contracting an $\Sc$ redex changes it by exactly $+\size{z}$. No function of circle count is therefore monotone along reductions, and the termination argument of the pure fragment does not transfer.
\item \emph{An explicit infinite reduction} (Section~\ref{sec:nontermination}). The term $\Wc\,\Wc$ with $\Wc=\Sc\,\Ic\,\Ic$ and $\Ic=\Sc\,\Kc\,\Kc$ has six circles in $\Wc$ and an infinite reduction. Under leftmost-outermost reduction the exact sequence is computed in closed form and the term grows without bound. Under three other strategies it cycles with period five.
\item \emph{Confluence, and its separation from strategy} (Section~\ref{sec:confluence}). Reduction is Church--Rosser, proved here by the parallel-reduction method, so a term has at most one normal form. This does not say that every strategy finds that normal form, and the term $\Kc\,a\,\Omega$ shows that it does not.
\item \emph{Exact translation from lambda terms} (Sections~\ref{sec:abstraction} and~\ref{sec:lambda}). Bracket abstraction satisfies $([x]M)\,N\reds M[x:=N]$ with a step count that is computed exactly from the structure of $M$. It follows that weak $\beta$ reduction is simulated step for step, and also that reduction under a binder is \emph{not} simulated, which is shown by an explicit pair of distinct normal forms.
\item \emph{Universality} (Section~\ref{sec:universality}). Counter machines compile to lambda terms, and thence to circles, in such a way that a machine halts with output $n$ exactly when the circle term reduces to the Scott numeral for $n$. The converse direction, that a non-halting machine yields a term with no normal form, depends on the leftmost normalization theorem, which is cited and not reproved.
\item \emph{A dynamic-programming reading} (Section~\ref{sec:dp}). Reduction graphs satisfy Bellman recursions for shortest and longest reduction and for path counts, the pure fragment's pop orderings are counted by the hook length formula and by an ideal-counting recursion, sharing is memoization, and the number of ways to group a flat expression is computed by an interval recurrence that sums to Catalan numbers.
\item \emph{Construction traces} (Section~\ref{sec:traces}). A reduction can be accompanied by a record of why a decomposition was chosen, and the obligations of that decomposition can be stated and checked exactly.
\end{enumerate}

\subsection{How the evidence is organized}

Each statement in the essay carries one of four statuses, and the status is part of the statement. A result is \emph{proved here} when a proof appears in the text. It is \emph{cited} when it is a published theorem whose proof is not reproduced, and the reference is given. It is \emph{machine-checked within stated bounds} when an executable companion verified it on an exhaustive enumeration or a stated random corpus, in which case the bound is part of the claim and the check is evidence and not proof. And it is \emph{not claimed} when it is something a reader might expect and the essay does not assert. Section~\ref{sec:ledger} collects these into a single table. The companion is a small program in JavaScript with an independent reference implementation in Python that shares no code and no data structures with it. All numbers appearing in the tables and in the running text are generated from the measured output of that program, so that no figure in the essay was typed by hand.

The scale of the checking is as follows. The reduction engine was exercised on every combinator word with up to six leaves over $\{\Sc,\Kc,x\}$, on every word with seven or eight leaves over $\{\Sc,\Kc\}$, and on a corpus of random words, and it was compared with the independent implementation on \factDiffComparisons{} strategy-by-term comparisons, of which \factDiffCycle{} ended in a cycle and \factDiffBudget{} and \factDiffSize{} in a step or size budget, so that the comparison is not confined to terminating cases. Bracket abstraction was verified, contraction by contraction, on \factAbsCases{} instances. The reader who distrusts any of this is invited to run the suite, to raise its bounds, and to read the oracle. The bounds are printed in the title of each test.

\subsection{Reading guide}

Sections~\ref{sec:pure} to~\ref{sec:size} establish the pure fragment, the extension, and the size laws. Section~\ref{sec:nontermination} is about nontermination and strategy. Section~\ref{sec:confluence} proves confluence. Sections~\ref{sec:abstraction} to~\ref{sec:universality} carry out the route to universality. Section~\ref{sec:lisp} states the relationship to Lisp, and Section~\ref{sec:dp} gives the dynamic-programming reading. Section~\ref{sec:traces} treats construction traces, and Section~\ref{sec:ledger} is the ledger.
